\documentclass[12pt]{article}

\usepackage{sbc-template}

\usepackage[T1]{fontenc}
\usepackage[utf8]{inputenc}
\usepackage[brazilian]{babel}
\usepackage{graphicx,url}
\usepackage{booktabs}
\usepackage{listings}
\usepackage{xcolor}

\lstset{
  basicstyle=\ttfamily\scriptsize,
  breaklines=true,
  breakatwhitespace=false,
  columns=fullflexible,
  frame=single,
  backgroundcolor=\color[gray]{0.95},
  xleftmargin=0pt,
  xrightmargin=0pt
}

\sloppy

\title{Autoescalamento Horizontal de Pods em Kubernetes: uma Análise
  Comparativa entre Implantações Local (Minikube) e em Nuvem (AWS EKS)}

\author{Ivo Motoki Makiyama Lopes\inst{1}}

\address{Escola Politécnica -- Universidade de São Paulo (USP)\\
  Disciplina PSI5120 -- Tópicos de Computação em Nuvem
  \email{ivo.motoki@usp.br}
}

\begin{document}

\maketitle

\begin{center}
\footnotesize Código-fonte, manifestos, scripts e evidências completas
disponíveis em: \url{https://github.com/IvoMotoki/psi5120-trabalho-hpa}
\end{center}

\begin{abstract}
  This paper presents the design, deployment, and evaluation of a web server
  running under Kubernetes Horizontal Pod Autoscaler (HPA), deployed in two
  environments: a local Minikube cluster and a managed AWS EKS cluster. Using
  an identical CPU-bound workload and load-generation method in both
  environments, we measured HPA reaction time, maximum replica count reached,
  scale-down time, approximate cost, and operational complexity. The Minikube
  deployment scaled from 1 to 10 replicas (the configured maximum) in about 2
  minutes and 10 seconds, unconstrained by node capacity. The EKS deployment,
  running on Free-Tier-eligible \texttt{t3.micro} nodes, reacted even faster
  (roughly 51 seconds to the first extra replica) but was capped at 5
  replicas by a per-node pod limit unrelated to CPU or memory, illustrating a
  structural difference between a resource-bound local cluster and a
  managed, quota-bound cloud environment.
\end{abstract}

\begin{resumo}
  Este artigo apresenta o projeto, a implantação e a avaliação de um servidor
  web sob autoescalamento horizontal de pods (HPA) no Kubernetes, implantado
  em dois ambientes: um cluster local via Minikube e um cluster gerenciado na
  AWS EKS. Usando a mesma carga de trabalho, intensiva em CPU, e o mesmo
  método de geração de carga nos dois ambientes, medimos o tempo de reação do
  HPA, o número máximo de réplicas atingido, o tempo de retração, o custo
  aproximado e a facilidade de gerenciamento. A implantação no Minikube
  escalou de 1 para 10 réplicas (o máximo configurado) em cerca de 2 minutos
  e 10 segundos, sem restrição de capacidade dos nós. A implantação no EKS,
  rodando em nós \texttt{t3.micro} elegíveis para o nível gratuito da AWS,
  reagiu ainda mais rápido (cerca de 51 segundos até a primeira réplica
  extra), mas ficou limitada a 5 réplicas por um limite de pods por nó
  independente de CPU ou memória, evidenciando uma diferença estrutural
  entre um cluster local limitado por recursos de hardware e um ambiente
  gerenciado na nuvem limitado por cotas de infraestrutura.
\end{resumo}

\section{Introdução}

O autoescalamento horizontal é um dos mecanismos centrais que tornam a
computação em nuvem atrativa para cargas de trabalho com demanda variável:
em vez de provisionar capacidade fixa para o pico de uso, o sistema ajusta
automaticamente o número de instâncias de um serviço em resposta a métricas
observadas, como utilização de CPU. No ecossistema Kubernetes, esse
mecanismo é implementado pelo \emph{Horizontal Pod Autoscaler} (HPA), um
controlador que monitora periodicamente métricas de recursos dos pods de um
\emph{Deployment} e ajusta o número de réplicas dentro de limites mínimo e
máximo configurados, de forma a manter a utilização média próxima de um
valor-alvo.

Embora o HPA seja conceitualmente simples, seu comportamento prático depende
fortemente do ambiente em que o cluster Kubernetes está implantado. Um
cluster local, como os providos pelo Minikube, tem sua capacidade
determinada apenas pelos recursos de CPU e memória alocados à máquina
virtual ou contêiner que hospeda o cluster. Já um cluster gerenciado em
nuvem, como o Amazon Elastic Kubernetes Service (EKS), introduz camadas
adicionais de infraestrutura -- rede virtual, balanceadores de carga,
integração com o provedor de identidade e acesso, e principalmente o tipo de
instância EC2 usado nos nós -- que podem impor restrições de capacidade
independentes da CPU ou memória efetivamente disponíveis.

Este trabalho tem como objetivo projetar, implantar e testar o mesmo
servidor web, sob a mesma política de HPA -- seguindo o roteiro oficial de
autoescalamento do Kubernetes~\cite{k8s-hpa-walkthrough} e a documentação de
uso do HPA no Amazon EKS~\cite{aws-eks-hpa} --, em duas implantações
distintas:
(A) um cluster Kubernetes local via Minikube e (B) um cluster gerenciado na
AWS EKS. A carga de trabalho e o método de geração de estresse foram
mantidos idênticos nas duas implantações, de modo a isolar as diferenças de
comportamento que decorrem do próprio ambiente de execução. A análise
compara o tempo de reação do HPA, o número máximo de réplicas efetivamente
atingido, o tempo de retração após o fim da carga, o custo aproximado de
cada implantação e a facilidade de gerenciamento observada durante o
processo.

O restante deste artigo está organizado da seguinte forma: a Seção~2 revisa
os conceitos de autoescalamento horizontal no Kubernetes; a Seção~3 descreve
a arquitetura e a topologia das duas implantações; a Seção~4 detalha a
metodologia de implantação e de geração de carga; a Seção~5 apresenta os
resultados obtidos e a análise comparativa; e a Seção~6 traz as conclusões e
considerações finais.

\section{Autoescalamento Horizontal de Pods no Kubernetes}

O HPA é um controlador nativo do Kubernetes que atua sobre um
\emph{Deployment} (ou outro objeto com escala, como um \emph{ReplicaSet} ou
\emph{StatefulSet}), ajustando o campo \texttt{replicas} periodicamente com
base em métricas coletadas de cada pod. No modo mais comum, a métrica usada
é a utilização de CPU, expressa como percentual do valor pedido em
\texttt{resources.requests.cpu} de cada contêiner: se a média de utilização
entre os pods prontos ultrapassa o percentual-alvo configurado, o
controlador calcula um número maior de réplicas necessário para trazer a
utilização média de volta ao alvo, respeitando o limite \texttt{maxReplicas};
se a utilização cai abaixo do alvo, o número de réplicas é reduzido,
respeitando o limite \texttt{minReplicas}.

Duas características do HPA são especialmente relevantes para os resultados
deste trabalho. Primeiro, o controlador depende de uma fonte externa de
métricas de recursos -- tipicamente o \emph{metrics-server}, um componente
que coleta periodicamente dados de uso de CPU e memória de cada nó/pod via
\texttt{kubelet} e os expõe pela API de métricas do Kubernetes
(\texttt{metrics.k8s.io}). Sem o metrics-server instalado e funcional, o HPA
não consegue calcular a utilização e permanece com o valor \texttt{<unknown>}
no lugar do percentual-alvo. Segundo, o HPA aplica políticas assimétricas de
estabilização: o aumento de réplicas (\emph{scale-up}) reage rapidamente
(por padrão, a cada ciclo de sincronização do controlador, tipicamente da
ordem de 15 segundos), enquanto a redução de réplicas (\emph{scale-down})
respeita, por padrão, uma janela de estabilização de 5 minutos contada a
partir da última recomendação de réplicas elevada -- uma medida para evitar
oscilações (\emph{flapping}) quando a carga é intermitente. Essa assimetria
é observada de forma consistente em ambas as implantações deste trabalho.

É importante distinguir o autoescalamento horizontal do autoescalamento
vertical (\emph{Vertical Pod Autoscaler}, VPA): enquanto o HPA aumenta ou
diminui o \emph{número} de réplicas de um pod, mantendo os requests/limits
de cada réplica fixos, o VPA ajusta os próprios requests/limits de CPU e
memória de um pod já em execução. Este trabalho usa exclusivamente o HPA,
por ser o mecanismo mais diretamente comparável entre uma implantação local
e uma implantação em nuvem gerenciada.

\section{Arquitetura e Topologia das Soluções}

\subsection{Implantação A -- Minikube (local)}

A Implantação A utiliza o Minikube (versão 1.38.1) para criar um cluster
Kubernetes de nó único, executado sobre o driver Docker, dentro de uma
distribuição WSL2 (Ubuntu 24.04) em uma máquina Windows 11. O perfil do
cluster (\texttt{hpa-2026}) foi criado com 3 vCPUs e 4096\,MB de memória
alocados à máquina virtual do Minikube. O \emph{metrics-server} foi
habilitado via o addon nativo do Minikube
(\texttt{minikube addons enable metrics-server}), que já aplica
internamente os ajustes necessários para funcionar em um ambiente de
desenvolvimento local (como a validação relaxada de certificados TLS do
kubelet). Toda a carga de trabalho foi implantada no \emph{namespace}
dedicado \texttt{tf-hpa}. Por se tratar de um cluster local de nó único, a
capacidade disponível para o autoescalamento é determinada exclusivamente
pela CPU e memória alocadas à máquina virtual -- não há um segundo nível de
quota de infraestrutura, como ocorre em um provedor de nuvem.

\subsection{Implantação B -- AWS EKS (nuvem)}

A Implantação B utiliza o Amazon EKS, um serviço gerenciado de Kubernetes,
provisionado via \texttt{eksctl} (versão 0.230.0) na região
\texttt{us-east-1}. O cluster (\texttt{psi5120-tf-hpa-eks}, Kubernetes
1.34) é composto por um \emph{control plane} totalmente gerenciado pela AWS
e um \emph{node group} gerenciado com instâncias EC2 do tipo
\texttt{t3.micro}. A escolha do tipo de instância não foi arbitrária: a
conta AWS utilizada está sujeita a uma restrição de elegibilidade para o
nível gratuito (\emph{Free Tier}), rejeitando a criação de instâncias
\texttt{t3.medium} originalmente planejadas com o erro
\texttt{InvalidParameterCombination: not eligible for Free Tier}. O node
group foi dimensionado com 4 nós -- um ajuste feito após observar que, com
apenas 2 nós, a capacidade de pods por nó do \texttt{t3.micro} (ver
Seção~5.2) era insuficiente até mesmo para agendar os próprios pods de
sistema do EKS, sem sobrar espaço para o servidor web sob teste.

O EKS instala automaticamente, como \emph{addons} gerenciados, os
componentes \texttt{vpc-cni} (rede de pods via ENIs da VPC), \texttt{kube-
proxy}, \texttt{coredns} e, relevante para este trabalho,
\texttt{metrics-server} -- diferentemente do Minikube, em que o
metrics-server precisa ser explicitamente habilitado como addon. O acesso
externo à aplicação foi provido por um \texttt{Service} do tipo
\texttt{LoadBalancer}, que a AWS provisiona automaticamente como um Classic
Load Balancer (CLB) com um nome DNS público, permitindo confirmar o
funcionamento da implantação a partir de fora do cluster.

\section{Metodologia}

\subsection{Aplicação Alvo e Manifestos}

Em ambas as implantações, o alvo do autoescalamento é o mesmo: um
\emph{Deployment} chamado \texttt{php-apache}, usando a imagem oficial de
exemplo do \emph{walkthrough} de HPA do Kubernetes
(\texttt{registry.k8s.io/hpa-example}) -- um servidor PHP/Apache com um
endpoint que executa um cálculo intensivo de CPU a cada requisição,
permitindo gerar carga de CPU de forma previsível e proporcional ao volume
de requisições. O contêiner é configurado com \texttt{requests.cpu = 200m} e
\texttt{limits.cpu = 500m} (\texttt{requests.memory = 64Mi},
\texttt{limits.memory = 128Mi}); o valor de \texttt{requests.cpu} é a base
sobre a qual o HPA calcula o percentual de utilização. O objeto
\texttt{HorizontalPodAutoscaler} associado usa a versão
\texttt{autoscaling/v2} da API, com \texttt{minReplicas: 1},
\texttt{maxReplicas: 10} e uma métrica de recurso do tipo \texttt{Utilization}
com alvo de 50\% sobre \texttt{requests.cpu}. Usar exatamente o mesmo par de
manifestos (Deployment e HPA) nas duas implantações é o que torna a
comparação entre elas metodologicamente válida: qualquer diferença de
comportamento observada decorre do ambiente de execução, não da carga de
trabalho em si.

A exposição do serviço difere ligeiramente entre as implantações por
necessidade prática: no Minikube, o \texttt{Service} é do tipo
\texttt{ClusterIP} (não há balanceador de carga de nuvem disponível
localmente); no EKS, um segundo \texttt{Service}, também do tipo
\texttt{LoadBalancer}, foi adicionado especificamente para demonstrar o
acesso público -- sem alterar o Deployment ou o HPA usados no teste de
carga em si.

\subsection{Geração de Carga}

A carga de estresse foi gerada por um conjunto de pods \texttt{busybox}
executando, em laço infinito, requisições HTTP (\texttt{wget}) contra o
\texttt{Service} \texttt{php-apache} pelo nome DNS interno do cluster
(\texttt{php-apache.tf-hpa.svc.cluster.local}). Gerar a carga de dentro do
próprio cluster, resolvendo o serviço pelo DNS interno, garante que o mesmo
mecanismo de teste seja usado nas duas implantações, independentemente da
topologia de rede externa de cada uma. O script de automação
(\texttt{scripts/load\_test.sh}) sobe os pods geradores, registra um
retrato do estado do HPA e dos pods a cada 10 segundos por um período de
carga ativa (240 segundos), remove os geradores de carga e continua
registrando por um período adicional (360 segundos) para observar a
retração. O número de pods geradores usado foi de 6 no Minikube e de 3 no
EKS -- a redução no EKS decorre diretamente da limitação de capacidade de
pods por nó discutida na Seção~5.2, e não de uma escolha metodológica
arbitrária.

\section{Resultados e Análise Comparativa}

\subsection{Implantação A -- Minikube}

A Tabela~\ref{tab:minikube} resume os principais momentos do ciclo de
autoescalamento observado no Minikube.

\begin{table}[ht]
\centering
\caption{Ciclo do HPA no Minikube}
\label{tab:minikube}
\begin{tabular}{lccc}
\toprule
Momento & Hora & Réplicas & CPU (alvo 50\%) \\
\midrule
Antes da carga            & 02:38:27 & 1  & 8\%   \\
Início do scale-up (CPU)  & 02:39:38 & 1  & 192\% \\
1\textsuperscript{o} scale-up          & 02:39:49 & 4  & 192\% \\
Scale-up intermediário    & 02:41:47 & 8  & 167\% \\
Pico (máximo configurado) & 02:41:59 & 10 & 167\% \\
Carga removida            & 02:42:35 & 10 & 131\% \\
CPU baixa, sem retração ainda & 02:44:28 & 10 & 10\% \\
Retração completa         & 02:54:40 & 1  & 26\%  \\
\bottomrule
\end{tabular}
\end{table}

O tempo entre a carga atingir os pods geradores (aproximadamente 02:38:39)
e o primeiro aumento de réplicas (02:39:49) foi de cerca de 1 minuto e 22
segundos. O HPA atingiu o número máximo de réplicas configurado
(\texttt{maxReplicas: 10}) em aproximadamente 2 minutos e 10 segundos desde
o início da carga, sem qualquer indício de limitação de capacidade -- a
única variável que limitou o escalonamento foi o próprio parâmetro
\texttt{maxReplicas}, não a CPU ou memória disponíveis na máquina virtual do
Minikube. Após a remoção da carga (02:42:35), a utilização de CPU caiu
rapidamente (10\% já às 02:44:28), mas o número de réplicas só retornou a 1
às 02:54:40 -- um intervalo de aproximadamente 12 minutos entre o fim da
carga e a retração completa, consistente com a janela padrão de
estabilização de \emph{scale-down} do HPA (5 minutos), contada a partir da
última leitura de utilização elevada, e não a partir do instante em que a
carga foi removida.

\subsection{Implantação B -- AWS EKS}

A Tabela~\ref{tab:eks} resume o ciclo equivalente observado no EKS.

\begin{table}[ht]
\centering
\caption{Ciclo do HPA no AWS EKS}
\label{tab:eks}
\begin{tabular}{lccc}
\toprule
Momento & Hora & Réplicas & CPU (alvo 50\%) \\
\midrule
Antes da carga                 & 01:06:57 & 1 & 8\%   \\
Início do scale-up (CPU)       & 01:07:31 & 1 & 116\% \\
1\textsuperscript{o} scale-up             & 01:07:48 & 3 & 248\% \\
Pico (limitado por capacidade) & 01:08:05 & 5 & 185\% \\
Carga removida                 & $\sim$01:10:57 & 5 & --    \\
CPU baixa, sem retração ainda  & 01:16:59 & 5 & 8\%   \\
Retração completa              & 01:18:05 & 1 & 10\%  \\
\bottomrule
\end{tabular}
\end{table}

No EKS, o tempo de reação foi ainda menor que no Minikube: cerca de 34
segundos entre a carga atingir os pods e o primeiro salto de utilização
(116\%), e cerca de 51 segundos até a primeira réplica extra ser criada
(01:07:48). Esse tempo de reação mais rápido é compatível com o fato de o
\texttt{metrics-server} do EKS já estar em execução como addon gerenciado
desde a criação do cluster, sem a janela inicial de "aquecimento" de coleta
de métricas observada no Minikube logo após a implantação.

O aspecto mais relevante da Implantação B, porém, não é a velocidade de
reação, e sim o \emph{teto} de escalonamento efetivamente atingido. Mesmo
com a utilização de CPU medida em 248\%/50\% e 185\%/50\% -- valores muito
acima do necessário para justificar réplicas adicionais dentro do limite
configurado de \texttt{maxReplicas: 10} -- o HPA não conseguiu escalar além
de 5 réplicas. A causa não é de CPU ou memória, e sim um limite estrutural
do tipo de instância: nós \texttt{t3.micro} suportam no máximo 4 pods por
nó, um limite imposto pelo número de endereços IP disponíveis por
\emph{Elastic Network Interface} (ENI) na implementação da VPC CNI da AWS,
e não pelos recursos computacionais do nó. Com os \emph{daemonsets} do
próprio sistema do EKS (\texttt{aws-node} e \texttt{kube-proxy}, um pod de
cada por nó) já ocupando parte dessa capacidade, e com os pods
\texttt{coredns} e \texttt{metrics-server} (2 réplicas cada, por padrão)
disputando o restante, a capacidade líquida disponível para o
\texttt{php-apache} e para os próprios pods geradores de carga, mesmo com o
node group ampliado para 4 nós, ficou restrita a poucas unidades -- valor
consistente com o teto de 5 réplicas observado.

Esse comportamento contrasta diretamente com o Minikube, onde a única
restrição de capacidade é a CPU/memória alocada à máquina virtual do
cluster, sem um segundo nível de cota de infraestrutura. Trata-se de uma
diferença estrutural entre um cluster local e um cluster gerenciado em
nuvem operando sob restrição de custo: o mesmo HPA, com a mesma política e
sob uma carga proporcionalmente maior, produz resultados de escalonamento
qualitativamente diferentes dependendo exclusivamente do tipo de instância
disponível.

Quanto à retração, a janela de observação do script de teste (360 segundos
após a remoção da carga) quase capturou o evento de \emph{scale-down}, mas
terminou minutos antes; a retração completa (5 para 1 réplica) foi
confirmada manualmente às 01:18:05, também consistente com a janela padrão
de estabilização de 5 minutos do HPA -- o mesmo comportamento assimétrico
observado no Minikube.

\subsection{Custo Aproximado}

A Implantação A tem custo direto nulo, por rodar inteiramente em recursos
de hardware já disponíveis na máquina do desenvolvedor. A Implantação B
incorre em custo real de infraestrutura AWS: a taxa fixa do control plane
gerenciado do EKS (aproximadamente US\$0,10/hora, sem cobertura do nível
gratuito) e o custo das instâncias EC2 do node group (\texttt{t3.micro},
tipicamente elegíveis ao nível gratuito da AWS dentro do limite de 750
horas/mês combinadas, ou aproximadamente US\$0,0104/hora cada, fora desse
limite). O cluster da Implantação B permaneceu ativo por aproximadamente 63
minutos, do provisionamento à exclusão completa confirmada, resultando em
um custo aproximado de US\$0,10 a US\$0,15 para todo o experimento -- a
maior parte proveniente do control plane, que não tem nível gratuito,
independentemente do tempo de uso.

\subsection{Facilidade de Gerenciamento}

A implantação local via Minikube foi operacionalmente mais simples de
iniciar -- um único comando cria o cluster e o addon de metrics-server em
poucos minutos --, mas exigiu atenção a recursos do host: uma tentativa
inicial de criação do cluster falhou de forma intermitente devido a
esgotamento de memória RAM da máquina Windows (menos de 500\,MB livres),
causando reinícios sucessivos do \emph{control plane} local até a memória
ser liberada. Esse tipo de instabilidade de ambiente é uma classe de
problema que praticamente não existe em um cluster gerenciado, onde o
control plane roda em infraestrutura da AWS, isolada da máquina do
desenvolvedor.

A implantação no EKS, por outro lado, exigiu etapas adicionais de
configuração ausentes no fluxo local: instalação de ferramentas de linha de
comando específicas (\texttt{aws-cli}, \texttt{eksctl}), configuração de
credenciais de acesso programático via IAM (distintas das credenciais de
login no console web), e um tempo de provisionamento consideravelmente
maior -- cerca de 15 a 20 minutos apenas para o \emph{control plane}, mais
alguns minutos para o \emph{node group}, contra menos de um minuto no
Minikube. A conta utilizada também impôs uma restrição não documentada de
antemão (elegibilidade a instâncias do nível gratuito), que só se
manifestou como um erro durante a criação do node group e exigiu ajuste do
tipo de instância e, por consequência, do dimensionamento do cluster. Por
fim, o encerramento correto de uma implantação em nuvem exige disciplina
adicional -- excluir o \texttt{Service} do tipo \texttt{LoadBalancer} antes
do cluster, para evitar que o balanceador fique órfão, e verificar
explicitamente, via linha de comando, que nenhum recurso cobrável
permaneceu ativo --, uma preocupação operacional inexistente no Minikube,
onde basta apagar o perfil local.

\section{Conclusão}

Este trabalho implantou e testou o mesmo servidor web sob a mesma política
de Horizontal Pod Autoscaler em duas implantações Kubernetes distintas --
Minikube local e AWS EKS gerenciado -- usando exatamente a mesma carga de
trabalho e o mesmo método de geração de estresse, de forma a isolar as
diferenças de comportamento decorrentes do próprio ambiente de execução.

Os resultados mostram que o mecanismo de autoescalamento reage de forma
qualitativamente semelhante nos dois ambientes -- aumento rápido de
réplicas (da ordem de dezenas de segundos a poucos minutos) seguido de uma
retração deliberadamente mais lenta (da ordem de vários minutos, ditada pela
janela de estabilização padrão do HPA) --, mas que o \emph{teto} efetivo de
escalonamento pode ser determinado por fatores completamente alheios à
política do HPA em si. No Minikube, esse teto foi o próprio
\texttt{maxReplicas} configurado, atingido sem qualquer restrição adicional.
No EKS, o teto real foi imposto por um limite de pods por nó associado ao
tipo de instância elegível ao nível gratuito da AWS -- uma restrição de
infraestrutura que não tem equivalente direto em um cluster local, e que só
se torna visível na prática, e não na documentação do HPA.

Do ponto de vista de facilidade de gerenciamento, o Minikube favorece a
iteração rápida e o baixo atrito inicial, ao custo de depender inteiramente
da estabilidade e da capacidade de hardware da máquina local. O EKS, por
sua vez, oferece um control plane gerenciado e isolado, mas introduz
complexidade operacional adicional -- configuração de credenciais, tempo de
provisionamento maior, restrições de conta e cuidados redobrados de
limpeza de recursos para controle de custo -- que devem ser considerados ao
planejar cargas de trabalho de produção sensíveis a custo e a capacidade de
escalonamento.

Como trabalhos futuros, seria valioso repetir o experimento em instâncias
EC2 de maior capacidade de rede (fora do nível gratuito), de forma a
isolar o efeito do limite de pods por nó observado neste trabalho, e
avaliar também o \emph{Cluster Autoscaler} do Kubernetes em conjunto com o
HPA, permitindo que o número de nós do cluster gerenciado se ajuste
automaticamente à demanda de pods, e não apenas o número de réplicas do
Deployment.

\bibliographystyle{sbc}
\bibliography{trabalho_hpa_psi5120_ivo_motoki}

\appendix

\section{Apêndice: Trechos das Evidências Brutas}

Esta seção reproduz trechos representativos dos arquivos de log brutos
gerados automaticamente durante os experimentos (comando/saída de
\texttt{kubectl get hpa} e \texttt{kubectl get pods}, capturados a cada 10
segundos pelo script \texttt{scripts/load\_test.sh}), correspondentes aos
momentos resumidos nas Tabelas~\ref{tab:minikube} e~\ref{tab:eks}. Os
arquivos completos (com todas as capturas intermediárias e o estado de cada
pod) estão disponíveis no repositório do projeto, nos diretórios
\texttt{evidencias/minikube/} e \texttt{evidencias/eks/}.

\subsection{Implantação A -- Minikube}

\textbf{Antes da carga (02:38:27):}
\begin{lstlisting}
NAME         REFERENCE               TARGETS       MINPODS   MAXPODS   REPLICAS   AGE
php-apache   Deployment/php-apache   cpu: 8%/50%   1         10        1          4m59s
\end{lstlisting}

\textbf{Primeiro scale-up (02:39:49):}
\begin{lstlisting}
NAME         REFERENCE               TARGETS         MINPODS   MAXPODS   REPLICAS   AGE
php-apache   Deployment/php-apache   cpu: 192%/50%   1         10        4          6m22s
\end{lstlisting}

\textbf{Pico -- máximo configurado atingido (02:41:59):}
\begin{lstlisting}
NAME         REFERENCE               TARGETS         MINPODS   MAXPODS   REPLICAS   AGE
php-apache   Deployment/php-apache   cpu: 167%/50%   1         10        10         8m31s
\end{lstlisting}

\textbf{Carga removida (02:42:35):}
\begin{lstlisting}
NAME         REFERENCE               TARGETS         MINPODS   MAXPODS   REPLICAS   AGE
php-apache   Deployment/php-apache   cpu: 131%/50%   1         10        10         9m8s
\end{lstlisting}

\textbf{Retração completa (02:54:40, confirmada manualmente após o fim da
janela de observação do script):}
\begin{lstlisting}
NAME         REFERENCE               TARGETS        MINPODS   MAXPODS   REPLICAS   AGE
php-apache   Deployment/php-apache   cpu: 26%/50%   1         10        1          21m
\end{lstlisting}

\subsection{Implantação B -- AWS EKS}

\textbf{Antes da carga (01:06:58):}
\begin{lstlisting}
NAME         REFERENCE               TARGETS       MINPODS   MAXPODS   REPLICAS   AGE
php-apache   Deployment/php-apache   cpu: 8%/50%   1         10        1          20m
\end{lstlisting}

\textbf{Primeiro scale-up (01:07:48):}
\begin{lstlisting}
NAME         REFERENCE               TARGETS         MINPODS   MAXPODS   REPLICAS   AGE
php-apache   Deployment/php-apache   cpu: 248%/50%   1         10        3          20m
\end{lstlisting}

\textbf{Pico -- limitado por capacidade de pods por nó, não por
\texttt{maxReplicas} (01:08:04), com réplicas do \texttt{php-apache} presas
em \texttt{Pending} por falta de capacidade agendável:}
\begin{lstlisting}
NAME         REFERENCE               TARGETS         MINPODS   MAXPODS   REPLICAS   AGE
php-apache   Deployment/php-apache   cpu: 185%/50%   1         10        5          21m
NAME                              READY   STATUS    RESTARTS   AGE     NODE
load-generator-7f8684cc8f-5q6sp   1/1     Running   0          69s     ip-192-168-39-22.ec2.internal
load-generator-7f8684cc8f-k7h7g   1/1     Running   0          69s     ip-192-168-28-221.ec2.internal
load-generator-7f8684cc8f-qxbbv   1/1     Running   0          69s     ip-192-168-28-221.ec2.internal
php-apache-56cf58f979-9sbx4       0/1     Pending   0          37s     <none>
php-apache-56cf58f979-lbwfj       0/1     Pending   0          22s     <none>
php-apache-56cf58f979-lzn77       1/1     Running   0          8m23s   ip-192-168-39-22.ec2.internal
\end{lstlisting}

\textbf{Retração completa (01:18:05, confirmada manualmente após o fim da
janela de observação do script):}
\begin{lstlisting}
NAME         REFERENCE               TARGETS        MINPODS   MAXPODS   REPLICAS   AGE
php-apache   Deployment/php-apache   cpu: 10%/50%   1         10        1          31m
\end{lstlisting}

\subsection{Evidência de Acesso Externo (EKS)}

Confirmação de que a aplicação implantada no EKS estava acessível
publicamente pela internet, via o DNS do Load Balancer provisionado pelo
\texttt{Service} do tipo \texttt{LoadBalancer}:

\begin{lstlisting}
$ curl -i http://a06f3a000408841fcbfd80005be8ce5f-992750954.us-east-1.elb.amazonaws.com/

HTTP/1.1 200 OK
Date: Wed, 26 Aug 2026 03:55:41 GMT
Server: Apache/2.4.10 (Debian) PHP/5.6.14
X-Powered-By: PHP/5.6.14
Content-Length: 3
Content-Type: text/html; charset=UTF-8

OK!
\end{lstlisting}

\end{document}
